\documentclass{amsart}
%\usepackage[english]{babel}

\usepackage{amsfonts}
\usepackage{amssymb}
\usepackage{color}

\newtheorem{thm}{Theorem}[section]
\newtheorem{cor}[thm]{Corollary}
\newtheorem{lem}[thm]{Lemma}
\newtheorem{prob}[thm]{Problem}
\newtheorem{prop}[thm]{Proposition}
\theoremstyle{definition}
\newtheorem{defn}{Definition}[section]
\newtheorem{exmp}{Example}[section]

\theoremstyle{remark}
\newtheorem{rem}{Remark}[section]
\newtheorem{notation}{Notation}
\renewcommand{\thenotation}{}
\numberwithin{equation}{section}
\newcommand{\interval}[1]{\mathinner{#1}}

\begin{document}

\title[On Runge type theorems]{On Runge type theorems for solutions to 
strongly uniformly parabolic operators }

\author{P.Yu. Vilkov}

\address[Pavel Vilkov]
{Siberian Federal University
                                                 \\
         pr. Svobodnyi 79
                                                 \\
         660041 Krasnoyarsk
                                                 \\
         Russia}
\email{pavel\_vilkov17@mail.ru}

\author{A.A. Shlapunov}

\address[Alexander Shlapunov]
{Siberian Federal University
                                                 \\
         pr. Svobodnyi 79
                                                 \\
         660041 Krasnoyarsk
                                                 \\
         Russia}
\email{ashlapunov@sfu-kras.ru} 

\address{Sirius Mathematics Center
Sirius University of Science and Technology, Olimpiyskiy ave. b.1, 354349 Sochi,
         Russia}
\email{shlapunov.aa@talantiuspeh.ru}

\subjclass {Primary 35K25; Secondary 35J60}
\keywords{strongly uniformly parabolic  operators, approximation theorems}

\begin{abstract}
Let $G_1, G_2 $ be  domains in ${\mathbb R}^{n+1}$, $n \geq 2$, such that $G_1 \subset G_2$ 
and the domain $G_1$ have rather regular boundary. We investigate the problem of 
approximation of solutions to strongly uniformly $2m$-parabolic  system $\mathcal L$ in the 
domain $G_1$  by solutions to the same system in the  domain 
$G_2$. First, we prove that the space  $S _{\mathcal L}(G_2)$ of solutions to the system 
$\mathcal L$ in the domain $G_2$  is dense in the space $S _{\mathcal L}(G_1)$, endowed with  
the standard Fr\'echet topology of the uniform convergence on compact subsets in $G_1$, if 
and only if  the complements $G_2 (t) \setminus G_1 (t)$ have no non-empty compact components 
in  $G_2 (t)$ for each $t\in \mathbb R$, where $G_j (t) = \{x \in {\mathbb R}^n: (x,t) \in 
G_j\}$. Next, under additional  assumptions on the regularity of the bounded domains $G_1$ and 
$G_1(t)$, we prove that solutions from the Lebesgue class $L^2(G_1)\cap S _{\mathcal L}(G_1)$ 
can be approximated by solutions from $S _{\mathcal L}(G_2)$ if and only if  the same 
assumption on the  complements $G_2 (t) \setminus G_1 (t)$, $t\in \mathbb R$, is fulfilled.
\end{abstract}

\maketitle

\section*{Introduction}
\label{s.Int}



Apparently, the Approximation Theory in Analysis begins with the famous Weierstra\ss{} 
theorem for continuous functions on segments of the real line, see \cite{W1885}, 
where polynomials were used as the approximating set. It appears that the matter 
has become significantly more complicated for complex functions of complex variable if one 
wants to approximate  them by the \textit{holomorphic} polynomials, see \cite{R1885}. Actually,  
C. Runge \cite{R1885} proposed to use the Cauchy kernel for the so-called rational 
approximation of holomorphic functions on open connected sets (plane domains) 
in $\mathbb C$ in the topology related to the uniform convergence on compact subsets of the 
domain. The notion of Runge's pair $\omega \subset \Omega$ of plane domains, for which the 
space ${\mathcal O}(\Omega) $ of holomorphic functions  in $\Omega$ is dense in the space 
$ {\mathcal O}(\omega) $ (endowed with the topology discussed above) gives rise to the 
investigation of topological/geometrical conditions on the pair. The analysis of continuous 
functions approximable on  compact sets  lead to the notion of the analytic capacity and the 
creation of the related theory, see, for instance, the paper by A.G. Vitushkin \cite{V67}. 
The matters were successfully extended to the theory of harmonic functions, see 
\cite{Merg56}, and even for solutions to general elliptic operators with the uniqueness 
condition in small, see for instance, \cite[ch. 4, 5]{Tark37}, where the bilateral fundamental 
solutions to the related operators were used instead of the Cauchy kernel. Of course, one 
should mention the theorem for (non-necessarily elliptic) operators with constant coefficients 
related to the uniform approximation by exponential-polynomial solutions in convex domains, 
see \cite{Mal56} (the elliptic case), \cite{Mal63}, \cite[Ch. VI]{Pal}.

As for the parabolic equations, the problem of the uniform approximation in the context 
of Runge's pairs was successfully solved for the heat equation in the papers \cite{J}, 
\cite{D80}  (see also \cite{GauTa10} for "the rational approximation" by the functions 
constructed with the use of the heat kernel).

On the other hand,  approximation theorems in various function spaces, where behaviour of the 
elements are controlled up to the boundary of the considered sets, appeared to be more 
important for applications, see, for instance,  the  V.P. Havin \cite{H68} for the analytic 
functions or the monograph \cite[Ch. 5--8]{Tark36} for the Sobolev solutions to systems of 
differential equations with surjective/injective symbols. For $L^2(G_1)$-solutions to the heat 
equation or to the parabolic Lam\'e type operator in a cylinder domain $G_1=\Omega \times (0,T) 
\subset {\mathbb R}^{n+1} $ we refer to recent papers \cite{ShHeat} and \cite{VKShLame}, respectively.

Naturally, one should explain why the similar approximation theorems were overlooked in 
the framework of general theory of parabolic equations. Actually, there are two principal 
reasons for this fact. The first one is the above mentioned theorem on the uniform 
approximation of solutions to operators with the constant coefficients in convex domains by the 
exponential-polynomial solutions covering many needs of approximation results in the Theory of 
Partial Differential Equations. The second reason is that typical problems for parabolic 
operators (such as the Cauchy problem in ${\mathbb R}^n \times [0,T)$ or the problems in the 
cylinder domains related with the Dirichlet or Neumann data on the lateral side, combined with 
the Cauchy data on the base of the domain) were handled mostly with the Fourier method of 
separation of variables in the Bochner-Sobolev spaces,  where approximation with respect to 
the functions of the space variables is crucial, see, for example, \cite{Lion69}, 
\cite{Tema79}, or by the integral representation method in the H\"older spaces where the 
approximation is not usually used, see, for instance, \cite{frid}, \cite{LadSoUr67}. 

On the other hand, the non-standard Cauchy problem for solutions to the parabolic equations 
in cylinder domains with the Cauchy data on a part of the lateral side of the cylinder 
behaves much more like the ill-posed Cauchy problem for solutions to elliptic systems, see, 
for instance, \cite{ShTaLMS}, \cite[Ch. 10]{Tark36} for the elliptic theory and \cite{KuSh}, 
\cite{MMT17}, \cite{VKShLame}, \cite{PuSh15}, for the parabolic theory in the Sobolev type 
spaces. The approximation theorems for Runge's pairs in the Lebesgue type spaces are crucial 
for this type of problems because they provide both dense solvability  and a  
possibility to construct the approximate solutions to the problems, see \cite{Merg56}, 
\cite{MaHa74}, \cite[Theorem 7.6]{ShTaLMS}, and  \cite[Corollary 3.5]{VKShLame}, 
respectively. Actually, the non-standard ill-posed Cauchy problem for the parabolic operators 
plays essential role in the development of non-invasive methods of Cardiology, see, for 
instance, \cite[\S 4]{KSU_ZAMM}. Please, note that formulations of such problems do not 
usually use initial data at a suitable time $t=t_0$ and hence we do not need to pose them in 
the cylinder domains.  

In the present paper we consider approximation theorems for solutions to the strongly uniformly 
parabolic matrix differential operator ${\mathcal L} = \partial _t + L$ on a strip ${\mathbb R}^{n} 
\times \mathcal{I} $, where $\mathcal{I} $ is an interval on the time axis, $L$ is a strongly 
elliptic operator with bounded regular coefficients on the strip admitting the bilateral 
fundamental solution and possessing the Unique Continuation Property with respect to the space 
variables. As far as the topic targets the problems without initial data, the presented 
approach is fit for both the parabolic and backwards parabolic operators.  

Section \S \ref{s.1} is devoted to the uniform approximation on compact subsets of the domain 
$G_1 \subset {\mathbb R}^{n+1}$ of the elements of the space $S_{\mathcal L} (G_1)$  
consisting of continuous solutions to the equation ${\mathcal L} u = 0$ in $G_1$ by the 
elements of the space $S_{\mathcal L} (G_2)$ where the domain $G_2$ contains $G_1$. Similarly 
to the case of the heat equation, the space $S_{\mathcal L} (G_2)$ is dense in $S_{\mathcal L} 
(G_1)$ if and only if for each $t \in {\mathcal I}$ the complement of the set $G_1 (t) = \{x \in 
{\mathbb R}^n: (x,t) \in G_1\}$  in the set $G_2 (t)$ has no compact (non-empty) components. 
Actually, taking in accounts a small gap in the proof of the approximation theorem for the  
heat equation in \cite{D80} (that was discovered by \cite{GauTa10}) concerned with rather 
general assumptions on the structure of the domains considered as Runge's pairs in this 
particular situation,  we additionally assume some regularity of the domains' boundaries, see 
assumption $(A)$, in section \S \ref{s.1}.

In the section \S \ref{s.2} we consider a more subtle problem of approximation 
of solutions to parabolic operator ${\mathcal L}$ from the Lebesgue class 
 $L^2(G_1)$  by more regular solutions in a bigger domain $G_2$. 
We present a solution to this problem in the case where $G_1$ is a bounded 
domain with piece-wise smooth boundary with additional geometric restrictions (see $(A1)$, 
$(A2)$). Finally, as a by-product, we obtain the theorem on existence of bases with double 
orthogonality property in Sobolev type spaces of solutions to parabolic systems.  

\section{The uniform approximation}
\label{s.1}

Let ${\mathbb R}^n$, $n \geq 1$, be the $n$-dimensional Euclidean space with the coordinates  
$x=(x_1, \dots , x_n)$ and let  $\Omega \subset {\mathbb R}^n$ be a bounded domain 
(open connected set). As usual, denote by  $\overline{\Omega}$ the closure of $\Omega$, and by
 $\partial\Omega$ its boundary.  
 
We consider functions over  ${\mathbb R}^n$ and 
${\mathbb R}^{n+1}$.  As usual, for $s \in {\mathbb Z}_+$ we denote by $C^s(\Omega)$ 
and $C^s(\overline \Omega)$  the spaces of all $s$ times continuously differentiable 
functions on $\Omega$ and $\overline \Omega$, respectively.  Next,  for  $p\in [0,1)$, 
we denote by $C^{s,p}(\overline \Omega)$ the standard H\"older spaces. 
The spaces $C^{s,p}(\overline \Omega)$ are known to be Banach spaces with the standard  
norms and the  $C^{s,p}(\Omega)$ are the Fr\'echet spaces with the standard semi-norms, see, 
for instance, \cite{Kry96}, \cite{LadSoUr67}.

Let  also $L^2 (\Omega)$ be the Lebesgue space over $\Omega$ with the standard inner product 
$(u,v)_{L^2 (\Omega)} $ and let $H^s (\Omega)$, $s\in \mathbb N$, be the Sobolev space  
with the standard inner product $(u,v)_{H^s (\Omega)}$. As usual, we 
consider the Sobolev space $H^{-s} (\Omega)$, $s\in \mathbb N$, as the dual space 
$H^{s}_0 (\Omega)$ where $H^{s}_0 (\Omega)$ is the closure of the space 
$C^\infty_{\rm comp} (\Omega)$ consisting of smooth 
functions with compact supports in $\Omega$.  

For a natural number $k$, it is convenient to denote by ${\mathbf C} ^{s,p} _k (\Omega)$ 
the space of $k$-vector functions with components of the class $C ^{s,p} (\Omega)$ and, 
similarly, for the space $\mathbf{L}^2_k (\Omega)$, etc. 

Let $\mathcal{I}$ be a finite or infinite interval on the time axis and let  $L$ be a 
$(k\times k)$-matrix differential operator 
with continuous coefficients in a strip ${\mathbb R}^{n} \times \mathcal{I}$ of an even order 
$2m$:
$$
L = \sum_{|\alpha|\leq 2m} L_{\alpha} (x,t) \partial ^\alpha_x 
$$ 
where $L_{\alpha} (x,t)$ are $(k\times k)$-matrices with entries as described above and 
such that $L^*_\alpha (x,t)= L_\alpha (x,t)$ for all multi-indexes $\alpha \in {\mathbb N}^n$ 
with $|\alpha|=2m$, all $t\in \mathcal{I}$   and all $x \in {\mathbb R}^n$. Consider  strongly 
uniformly (Petrovsky) $2m$-parabolic operator  
$$
{\mathcal L}= \partial _t + L,
$$  
see, for instance, \cite{eid}, \cite{sol}. More precisely, this means that 
there is a positive constant $c_0$ such that 
$$
(-1)^{m} w^* \Big(
\sum_{|\alpha| =2m}L_\alpha (x,t)\zeta ^\alpha \Big) w \geq c_0 |w|^2 |\zeta|^{2mk}
$$
for all $t\in \mathcal{I}$, all $x \in {\mathbb R}^n$, all $\zeta \in {\mathbb R}^n \setminus 
\{0\}$ and all $w \in {\mathbb C}^k \setminus \{0\}$; here $w^*$ is the transposed and complex 
adjoint vector for the complex vector $w \in {\mathbb C}^k$. In particular, for each fixed 
$t_0\in \mathcal{I}$ the operator $L$ is (Petrovsky) elliptic with respect to the space 
variables $x$, i.e. 
$$
\det \Big(\sum_{|\alpha| =2m}L_\alpha (x,t_0)\zeta ^\alpha \Big) \ne 0
$$
for all $x \in {\mathbb R}^n$ and all $\zeta \in {\mathbb R}^n \setminus \{0\}$. 

Following \cite[pp. 245-246]{eid}, we assume also that the operator $\mathcal L$ satisfies 
the following assumptions:
\begin{enumerate}
\item[($\alpha_1$)] 
the coefficients $L_\alpha(x,t)$ are uniformly continuous over the strip 
${\mathbb R}^{n} \times \mathcal{I}$ for $|\alpha|=2m$; 
\item[($\alpha_2$)] 
the coefficients $L_\alpha(x,t)$ are uniformly bounded over the strip 
${\mathbb R}^{n} \times \mathcal{I}$; 
\item[($\alpha_3$)] 
the coefficients $L_\alpha(x,t)$ satisfy the H\"older condition  with respect to the space 
variables $x$ uniformly over the strip ${\mathbb R}^{n} \times \mathcal{I}$; 
\item[($\alpha_4$)]
each coefficient $L_\alpha(x,t)$ has partial derivative $\partial^\alpha _x L_\alpha(x,t)$ 
the strip ${\mathbb R}^{n} \times \mathcal{I}$ satisfying assumptions $(\alpha_1)$, 
$(\alpha_2)$, $(\alpha_3)$.
\end{enumerate}
Under these assumptions on  the coefficients of the operator it admits a unique fundamental 
solution  $\Phi (x,y,t,\tau)$ possessing standard estimates \cite[formulas (2.16), 
(2.17)]{eid}) and the normality property (\cite[Property 2.2]{eid}), i.e. 
\begin{equation} \label{eq.right}
{\mathcal L}_{x,t} \Phi(x,y,t,\tau) = I_k \delta (x-y, t-\tau),   
\end{equation}
where $\delta (x, t)$ is the Dirac functional supported at the point $(x,t)$ and  
\begin{equation} \label{eq.left}
{\mathcal L}^*_{y,\tau} \Phi^*(x,y,t,\tau)  =I_k \delta (x-y, t-\tau). 
\end{equation}

We are going to investigate solutions to $2m$-parabolic equations in  non-cylinder domains 
of special type, see for instance \cite[\S 22]{sol}. Namely let  $G$ be a  
domain in the strip ${\mathbb R}^{n} \times  \mathcal{I}$ and let 
$$
T_1(G)=T_1=\inf_{(x,t)\in G} t, \, T_2(G) = T_2=\sup_{(x,t)\in G} t. 
$$

Consider bounded sets $G (t) = \{x \in {\mathbb R}^n: (x,t) \in G\}$, $t\in 
\mathcal{I}$, playing an essential role in the sequel. We assume that the boundary 
of $G$ satisfy the following property.

\begin{itemize}
\item[(A)] 
For any numbers $t_3, t_4$ such that $T_1<t_3<t_4<T_2$ 
the set $\Gamma _{t_3,t_4}=\cup_{t\in [t_3,t_4]} \partial G (t)$ is a Lipschitz surface 
in ${\mathbb R}^{n+1}$. % if $n\geq 2$  or the union of finite numbers of intervals if $n=1$.
\end{itemize}

For us, the primary interest for studying parabolic equations in such domains 
is a possibility of applications in the Cardiology, where $G(t) \subset 
{\mathbb R}^3 \times (T_1,T_2)$ is the shape of the human 
myocardium at the time $t$ (see, for instance, \cite{Beheshti2016} for the related 
mathematical models or \cite{KSU_ZAMM} for the particular bi-domain model). 

With this purpose  we need the standard Banach anisotropic spaces $C^{2sm,s} (\overline G)$, see for instance, \cite[\S 22]{sol}, with the norm: 
$$
\|u\|_{C^{2ms,s} (\overline G)} = \sum_{|\alpha|+2mj\leq 2ms} 
\|\partial^\alpha_x \partial^j_t  u\|_{C (\overline G) }.
$$ 
The corresponding anisotropic spaces $C^{2sm,s} (G)$ are Fr\'echet spaces with the topology 
of the uniform convergence on compact subsets of $G$ with all the derivatives $\partial_x
^\alpha \partial _t ^j$, $|\alpha|+ 2mj\leq 2ms$, see, for instance, \cite{Sch71}.

Now, $S _{\mathcal L}(G)$ be the set of all the %generalized 
continuous $k$-vector functions on $G$, satisfying (homogeneous)  equation
\begin{equation} \label{eq.heat}
{\mathcal L} u = 0 \mbox{ in } G 
\end{equation}
in the sense of distributions. We endow this space  with the standard  topology of the 
uniform convergence on compact subsets of $G$. Next, we note that estimates  
\cite[formulas (2.16), (2.17)]{eid}) for the fundamental solution imply the standard interior 
a priori estimates for solutions of the operator $\mathcal L$, see, for instance, 
\cite[\S 19]{sol}, or \cite[Ch. 4, \S 2]{frid} for the second order operators. This means 
that all the continuous solutions to equation \eqref{eq.heat} are $(2m,1)$-differentiable on 
their domain, i.e. the following continuous embedding holds true:
$$
S _{\mathcal L}(G) \subset \mathbf{C}^{2m,1} _k (G).
$$ 
In particular, this means that $S _{\mathcal L}(G)$ is a closed subspace 
in $\mathbf{C} _k (G)$ and it is a Fr\'echet space itself (with the 
standard Fr\'echet topology inducing the standard uniform convergence together with all the 
derivatives on compact subsets of $G$). 

We also need one more assumption on the operator $\mathcal L$: the Unique Continuation 
Property with respect to the space variables. Namely, 
\begin{enumerate}
\item[$(UCP)$] For any solution $u \in S _{\mathcal L}(G)$ and any 
$t_0 \in (T_1,T_2)$, if $u (x,t_0)= 0$  for all $x$ from an open 
 subset $\omega \subset G (t_0)$ then $u \equiv 0$ in the open connected component 
of $G (t_0)$, containing $\omega$.
\end{enumerate}

Of course, if the ${\mathcal L}$ is the operator with constant coefficients then all the 
assumptions $(\alpha_1)$,   $(\alpha_2)$,  $(\alpha_3)$,  $(\alpha_4)$ and $(UCP)$ are 
fulfilled.  Besides, they hold true if  coefficients of the operator $\mathcal L$ are smooth, 
bounded and real analytic with respect to the variables $x$ for each $t\in \mathcal{I}$. In 
this case, the elements of $S_{\mathcal L}(G) $ are actually smooth in $G$ and they are real 
analytic with respect to the space variable $x\in G(t)$ for all $t \in (T_1,T_2)$, see, for 
instance \cite{eid}. Moreover, the following continuous embedding 
holds true in this particular situation:
$$
S _{\mathcal L}(G) \subset \mathbf{C}^{\infty} _k (G).
$$
  
Next, following C. Runge \cite{R1885} and \cite{D80}, we call domains 
$G_1 \subset G_2 \subset {\mathbb R}^{n+1}$ the ${\mathcal L}$-Runge's pair 
 if $S_{\mathcal L}(G_2)$ is everywhere dense in $S_{\mathcal L}(G_1)$. 
The following approximation theorem is quite similar to the corresponding statement 
for the heat equation, see \cite{J} for the case $G_2 = {\mathbb R}^{n+1}$ 
or \cite{D80} for the general case.  

\begin{thm} \label{t.dense.uniform}
Let ${\mathcal L}$ satisfy assumptions $(\alpha_1)$, $(\alpha_2)$, $(\alpha_3)$, 
$(\alpha_4)$, $(UCP)$. Let  also $G_1 \subset G_2 $   be domains 
in the strip ${\mathbb R}^n \times {\mathcal I}$ and $G_1$ 
satisfy assumption $(A)$. Then $S_{\mathcal L}(G_2)$ is everywhere dense in 
$S_{\mathcal L}(G_1)$ if and only if for each $t \in {\mathcal I}$ the complement of the 
set  $G_1 (t) $ has no compact (non-empty) components in the set $G_2 (t)$. 
\end{thm}

\begin{proof} We slightly modify  the proof from \cite{D80} for the solutions  
to the heat equation, using the duality theorems from modern functional analysis, 
see, for instance, \cite{Edw65} and the classical a priori estimates for parabolic equations, 
see, for example, \cite{eid}, \cite{frid}.  

As we have noted above, the space $S_{\mathcal L}(G_1)$ is a closed subspace of the space 
${\mathbf C}_k (G_1)$, endowed with the standard topology of the uniform convergence 
on compact subsets of  $G_1$. Then the Khan-Banach Theorem implies that $G_1$, $G_2$ is a 
${\mathcal L}$-Runge's pair if and only if any continuous functional $F$ on ${\mathbf C}_k 
(G_1)$ annihilating the space $S_{\mathcal L}(G_2)$ also annihilates the space 
$S_{\mathcal L}(G_1)$.

On the other hand, according the Riesz Theorem, see, for instance, 
\cite[Theorem 4.10.1]{Edw65}, any element $F$ of the dual space ${\mathbf 
C}_k^* (G_1)$ for ${\mathbf C}_k (G_1)$
can be presented with the use of a ($k$-vector valued) Radon measure $\mu_F$ with 
the compact support $K(\mu_F)$ in $G_1$ , i.e.
\begin{equation} \label{eq.Riesz}
F(u) = %\int_{K(\mu)} u(x,t) ,d\mu_F (x,t) 
\langle u ,d\mu_F  
\rangle \mbox{ for all } u \in S_{\mathcal L}(G_1).
\end{equation}

Of course, the operator ${\mathcal L}^*$ is backwards-parabolic and, for any 
solution $v(y,\tau)$ to the equation ${\mathcal L}^* v =0$, the vector 
$w(y,\tau)=v(y,-\tau)$  is a solution to parabolic equation 
$(\partial_\tau + L^* _y) w =0$. Thus, if $F\in {\mathbf C}_k^* (G_1)$ annihilates the space $S_{\mathcal L}(G_2)$ then, according to  
\eqref{eq.right}, for the following vector $W$ we have  
\begin{equation}
\label{eq.v1.A}
W(y,\tau)= %\int_{K(\mu)} 
\langle \Phi^* (x,y, t,\tau) , d\mu _F (x,t) \rangle=  0 \mbox{ for all } 
(y,\tau) \not \in G_2.
\end{equation}
But \eqref{eq.left} implies that 
\begin{equation} \label{eq.L*W.1}
{\mathcal L}^* W = d\mu_F \mbox{ in } {\mathbb R}^{n} \times \mathcal{I},  
\end{equation}
in the sense of distributions and, in particular, 
\begin{equation} \label{eq.L*W.2}
{\mathcal L}^* W = 0 \mbox{ in } %{\mathbb R}^{n+1}
\Big({\mathbb R}^{n} \times \mathcal{I} \Big) \setminus K(\mu_F).
\end{equation}
Therefore $W(y,\tau)$ is $C^{2m}$-smooth with respect to  $y$ in ${\mathbb R}^n 
\setminus K(\mu_F) (\tau)$ for each $\tau \in \mathcal{I}$ where 
$K(\mu_F) (\tau) = \{x\in {\mathbb R}^n: (x,\tau) \in K(\mu_F)\}$. 

However the complement $G_2 (t) \setminus G_1 (t)$ has no compact components in $G_2 (t)$ and 
hence each component of  ${\mathbb R}^{n} \setminus G_1 (t) $ intersects with
 ${\mathbb R}^{n} \setminus  G_2 (t) $ by a non-empty open set for each $t\in (T_1,T_2)$. 
Thus, \eqref{eq.v1.A} and the Unique Continuation Property $(UCP)$ imply that 
\begin{equation*} 
W (y,\tau) = 0   \mbox{ in }{\mathbb R}^{n} \setminus \overline{G_1 (\tau)}
\mbox{ for all } \tau\in \mathcal{I}, 
\end{equation*}
i.e. the vector $W$  is supported in $\overline G_1$. Following \cite{J}, we may even say 
that $W$ is supported in 
$$
\cup_{\tau \in (T_1,T_2) } \hat K_{G_1}(\mu_F) (\tau)
\subset \overline G_1, 
$$  
where $\hat K_{G_1}(\mu_F) (\tau)$ is the union of $K(\mu_F) (\tau)$ with  all the the 
components of $G_1 (\tau)\setminus K(\mu_F) (\tau)$ that are relatively compact in 
$G_1 (\tau)$.

Of course, if $W$ is compactly supported in $G_1$  
then we may take  a function $\varphi \in C^\infty (G_2) $ compactly supported in 
$G_1$ such that $\varphi \equiv 1 $ on a neighbourhood of ${\rm supp} \, (W)$. 
Then, by the Leibniz rule, for each $u \in S_{\mathcal L}(G_1)$ we have 
$$
{\mathcal L} (\varphi u) = \varphi {\mathcal L}  u   + 
(\partial_t \varphi) u  + \sum_{ |\alpha|\leq 2m} L_\alpha (x) 
\sum_{\beta + \gamma =\alpha , |\gamma|\geq 1} C_{\beta,\gamma}
(\partial ^{\gamma} \varphi) (\partial ^{\beta} u) = %\psi \tilde L u
$$
$$
(\partial_t \varphi) u  + \sum_{ |\alpha|\leq 2m} L_\alpha (x) 
\sum_{\beta + \gamma =\alpha , |\gamma|\geq 1} C_{\beta,\gamma}
(\partial ^{\gamma} \varphi) (\partial ^{\beta} u), 
$$
where $C_{\beta,\gamma}$ are binomial type coefficients; in particular,  
$$
{\rm supp} \, ({\mathcal L} (\varphi u)) 
\cap {\rm supp} \, (W)  =\emptyset.
$$ 
Therefore \eqref{eq.Riesz} and \eqref{eq.L*W.1} yield
\begin{equation} \label{eq.supported}
F (u) = \langle u ,d\mu_F  \rangle = \langle  u , {\mathcal L}^* W \rangle =
\langle  (\varphi u) , {\mathcal L}^* W \rangle = \langle  {\mathcal L} (\varphi u) ,  W 
\rangle = 0,
\end{equation} 
i.e. $F$ annihilates $S_{\mathcal L}(G_1)$, too. 

As Diaz \cite[p. 644]{D80} noted, complications may arise if  $W$  is not compactly supported 
in $G_1$  even in the case where ${\mathcal L}$ is the heat operator. Nevertheless  he 
considered more general type of domains and for this reason the related proof was rather 
complicated. 

In our more simple situation  we see that, as $\mu_F$ is compactly supported in $ G_1$, there 
are numbers $T_3$, $T_4$ such that $T_1<T_3<T_4<T_2$ and $K(\mu_F) (\tau) =\emptyset $ for 
all $\tau <T_3$ and $\tau>T_4$, i.e. $W$ is supported in $\cup_{\tau \in [T_3,T_4] } \hat 
K_{G_1}(\mu_F) (\tau)$. Moreover by the the assumptions $(\alpha_1)-(\alpha_4)$ on the 
operator ${\mathcal L}$, the vector $W$ belongs to ${\mathbf C}^{2m,1}_k 
(\Big({\mathbb R}^{n}\times \mathcal{I}\Big)\setminus K(\mu_F))$ and, 
in particular, $W=0$ in ${\mathbf C}^{2m,1}_k (\Big({\mathbb R}^{n}\times \mathcal{I}\Big)
\setminus G_1)$.

Now, according to assumption $(A)$, the union $\Gamma _{T_3,T_4}=\cup_{t\in [T_3,T_4]} 
\partial G (t)$ is a Lipschitz surface in ${\mathbb R}^{n+1}$. If necessary, 
increasing $T_4$ and decreasing $T_3$ we see that there is a domain 
$\tilde G \subset G_1$ with Lipschitz boundary consisting of the surface 
$\Gamma _{T_3,T_4}$ and the sets $G(T_3)$ and $G(T_4)$ such that the measure 
$\mu_F$ is compactly supported in $\tilde G$. 
Let $\tilde G_\varepsilon$ be the sets in $\tilde G$ such that the distance between 
the boundaries of $\tilde G_\varepsilon$ and $\tilde G$ be less than $\varepsilon$. 
Then there is a number $\varepsilon_0>0$ such that 
\begin{itemize}
\item
the set $\tilde G_{\varepsilon_0}$ is again is a domain with Lipschitz boundary and  the 
measure $\mu_F$ is compactly supported in $\tilde G_{\varepsilon}$. 
\item 
$\tilde G \setminus \tilde G_{\varepsilon_0}$ has no compact components in $\tilde G$.
\end{itemize} 
Replacing the support of measure $\mu_F$ with $\overline {\tilde G_{\varepsilon_0}}$, we see 
that $\cup_{\tau \in [T_3,T_4] } \hat K_{G_1}(\mu_F) (\tau) \subset \overline 
{\tilde G_{\varepsilon_0}}$ and hence $W$ is compactly supported in both $\tilde G$ and $G_1$.

Back, let there be a number $t_0\in \mathbb R$ such that the set  $G_1 (t_0) $   
have a compact (non-empty) component $K_0$ in the set $G_2 (t_0)$. 
in ${\mathbb R}^n$. Fix a point $y_0 \in K_0$. 
Then any vector column $U_j (x,t)$, $1\leq j \leq k$, 
of the fundamental matrix $\Phi (x-y_0, t-t_0)$ 
belongs to the space $S_{\mathcal L} (G_1)$. If the vector function 
$U_j (x,t)$ can be approximated in ${\mathbf C}^{2m,1}_k  (G_1)$ 
by a sequence $\{ u^{(i)} \}_{i \mathbb N}$ from the space $S_{\mathcal L} (G_2)$ 
then this sequence should satisfy the Cauchy criterion in the metrizable space 
${\mathbf C}^{2m,1}_k  (G_1 \setminus G_3) $ with a neighbourhood $G_3$ of the 
point $(y_0, t_0)$. Then by the standard a priori estimates for parabolic systems we conclude 
that sequence $\{ u^{(i)} \}_{i \mathbb N}$ converges, actually, to a vector 
function $U(x,t)$ in the space of smooth $k$-vector functions over a neighbourhood 
of the point $(y_0,t_0)$. In particular, assumption $(UCP)$ for the system ${\mathcal L}$ 
with respect to the space variables implies that the vector function $U_j (x,t) \in 
S_{\mathcal L} (({\mathbb R}^n \times \mathcal{I})\setminus (y_0,t_0))$ uniquely extends to a 
neighbourhood $G_4$ of the point $(y_0,t_0)$ as a solution $\tilde U_j (x,t)$ to 
${\mathcal L}$. Thus, we obtain a contradiction because for the matrix $\tilde \Phi (x,t)$
with columns $\tilde U_j (x,t)$, $1\leq j \leq k$, we have ${\mathcal L}\tilde \Phi (x,t) =0$ 
in $G_4$   but ${\mathcal L} \Phi  (x-y_0,t-t_0)$ coincides with the $\delta$-functional 
concentrated at the point $(y_0,t_0)$. 
\end{proof}

\section{The approximation in the mean}
\label{s.2}

In this section we discuss an approximation theorem for solutions to the operator 
$\mathcal L$ belonging to the Lebesgue spaces. Actually, it is quite similar to the 
approximation theorems for elliptic operators mentioned in the introduction and 
the approximation theorem with uniform convergence on compact subsets for parabolic systems 
proved in the previous section. Also, they are known for the heat equation or for the 
parabolic Lam\'e system in cylinder domains with rather regular lateral surfaces, 
see, for instance,  \cite{ShHeat} or \cite{VKShLame}. 

Investigating spaces of solutions to $2m$-parabolic equation, we need the 
anisotro\-pic Sobolev 
spaces than $H^{2ms,s} (G)$, $s \in  {\mathbb Z}_+$, in a domain $G\subset {\mathbb R}^{n} 
\times \mathcal{I}$ with the standard inner product,
$$
(u,v)_{H^{2ms,s} (G)} = \sum_{|\alpha|+2mj\leq 2ms} 
(\partial^\alpha_x \partial^j_t  u, \partial^\alpha_x \partial^j_t v)_{L^2 (G) }
$$ 
Also, for $\gamma \in {\mathbb Z}_+$, we denote by $H^{\gamma,2sm,s} (G)$ 
the set of all functions  $u \in H^{2sm,s} (G)$ such that 
 $\partial ^\beta_x u \in H^{2ms,s} (G)$ for all  $|\beta|\leq \gamma$.  
As before, it is convenient to denote by 
$\mathbf{H}^{2ms,s}_k (G)$ the space of all the $k$-vector functions with 
the components from ${H}^{2ms,s} (G)$, and similarly for the spaces
$\mathbf{H}^{\gamma, 2ms,s}_k (G)$, etc. 


We also will use the so-called Bocher spaces 
of functions depending on $(x,t)$ from the strip  
$\mathbb{R}^n \times  [T_1,T_2]$ with finite numbers $T_1<T_2$. 
Namely, for a Banach space $\mathcal B$ (for example, the space of functions 
on a sub-domain of $\mathbb{R}^n$) and    $p \geq 1$, we denote by  
$L^p (I,{\mathcal B})$ the Banach space of all the measurable mappings 
  $u : [T_1,T_2] \to {\mathcal B}$
with the finite norm  
$$
   \| u \|_{L^p ([T_1,T_2],{\mathcal B})}
 := \| \|  u (\cdot,t) \|_{\mathcal B} \|_{L^p ([T_1,T_2])},
$$
see, for instance, \cite[ch. \S 1.2]{Lion69},  \cite[ch.~III, \S~1]{Tema79}. 

The space $C ([T_1,T_2],{\mathcal B})$ is introduced with the use of the 
same scheme; this is the Banach space of all the continuous mappings
$u : [T_1,T_2] \to {\mathcal B}$ with the finite norm 
$$
   \| u \|_{C ([T_1,T_2],{\mathcal B})}
 := \sup_{t \in [T_1,T_2]} \| u (\cdot,t) \|_{\mathcal B}.
$$

Let ${\mathbf H}^{\gamma,2sm,s} _{k,\mathcal L}(G) =
\mathbf{H}^{\gamma,2sm,s}_k (G) \cap S _{\mathcal L}(G) $, 
$s \in  {\mathbb Z}_+$, $\gamma \in {\mathbb Z}_+$. 
By the discussion in Section \S \ref{s.1}, the space
${\mathbf H}^{2m,1} _{k,\mathcal L}(G)$ is a closed subspace of the Sobolev space 
$\mathbf{H}^{2m,1}_k(G)$. 
Similarly, if  coefficients of the operator $\mathcal L$ are smooth, 
bounded and real analytic with respect to the variables $x$ for each $t\in {\mathcal I}$,  
then ${\mathbf H}^{\gamma,2sm,s} _{k,\mathcal L}(G)$, ${\mathbf C}^{2ms,s} _{k,\mathcal L} 
(\overline{G}) = \mathbf{C}_k^{2ms,s} (\overline{G}) \cap S _{\mathcal L}(G) $, 
${\mathbf C}^{\infty} _{k,\mathcal L} (\overline{G}) =
\mathbf{C}_k^{\infty} (\overline{G}) \cap S _{\mathcal L}(G) $ 
are closed subspaces, consisting of solutions to equation \eqref{eq.heat}, in the spaces 
${\mathbf H}^{\gamma,2sm,s} _{k}(G)$, $\mathbf{C}^{2ms,s}_k  (\overline{G})$ and  
$\mathbf{C}^{\infty}_k  (\overline{G})$, respectively. 

Also, we need the space  $S _{\mathcal L}(\overline{G_2})$, 
defined as follows: 
$$
\cup_{G' \supset 
\overline{G} } S _{\mathcal L}(G'),
$$
where the union is with respect to all the domains $G' \subset {\mathbb R}^{n} \times 
(t_1, t_2)$,  containing the closure of the domain  $G$.
It follows from the a priori estimates discussed in \S \ref{s.1}
that the following (continuous) embeddings 
\begin{equation} \label{eq.emb.H}
S _{\mathcal L}(\overline{G}) \subset 
\mathbf{C}^{2m,1} _{k,\mathcal L} (\overline{G}) \subset
\mathbf{H}^{2m,1} _{k,\mathcal L}(G) 
\end{equation}
are fulfilled. Of course, if we additionally know that the coefficients of the operator 
${\mathcal L}$ are constant or smooth and real analytic with respect to the space variables 
then the operator is hypoelliptic and the following (continuous) embeddings 
\begin{equation} \label{eq.emb.H.infty}
S _{\mathcal L}(\overline{G}) \subset 
\mathbf{C}^{\infty} _{k,\mathcal L} (\overline{G}) \subset
\mathbf{H}^{\gamma,2ms,s} _{k,\mathcal L}(G_2) 
\end{equation} 
hold true  $\gamma,s\in {\mathbb Z}_+$. 

To prove an approximation theorem for the spaces of the Lebesgue solutions to 
$\mathcal L$,  we need more regularity of $\partial G$:  
\begin{itemize}
\item[(A1)]    
For each $t\in [T_1, T_2]$, the sets  
$G (t) = \{x \in {\mathbb R}^n: (x,t) \in G\}$, 
are  domains in ${\mathbb R}^n$ with $C^{2m}$-boundaries if $n\geq 2$ 
 or the union of finite numbers of intervals if $n=1$.
\item[(A2)] 
The boundary $\partial G$ of is the union $G (T_1)\cup  G (T_2) \cup \Gamma $, where 
$$\Gamma =\cup_{t\in (T_1,T_2)} \partial G (t)$$ 
is a $C^{2m,1}$-smooth surface of with 
no points where the tangential planes are parallel to the coordinate plane $\{t=0\}$, i.e. 
we have 
$$
\sum_{j=1}^n  (\nu_j (x,t))^2 \geq \varepsilon _0 \mbox{ for all } (x,t) \in \Gamma
$$
with a positive number $\varepsilon _0$.
\end{itemize}

Under these assumptions we easily see that functions from the space $H^{2m,1} (G)$ and  
some of their partial derivatives have reasonable traces on $\partial G$. The definition, 
the uniqueness and an existence theorem for traces of functions from anisotropic spaces can 
be found in \cite[\S 10]{BeINi} (see also \cite[Ch 2.]{UDP} for non-regular cases $n=2$). In our particular situation we may 
specify the traces  by hands (cf. \cite[Ch.3, \S 7]{MikhX} for anisotropic spaces in cylinder 
domains). With this purpose, we denote by $Y_{s-1/2} (\Gamma)$ the closure  of ${C}^{s} 
(\Gamma)$, $s \in \mathbb N$, with respect to the norm 
$$
\|\cdot\|_{Y_{s-1/2}(\Gamma)} = \int_{T_1}^{T_2} 
\|\cdot\|_{H^{s-1/2}_k (\partial G (t))} dt;
$$
by the construction, these are Hilbert spaces embedded continuously 
into ${L} ^2 (\Gamma)$. 

Of course, if $G = \Omega \times (T_1,T_2)$ is a cylinder domain in ${\mathbb R}^{n+1}$
with the base $\Omega$ being a domain with the boundary of class $C^{s}$, then 
$\Gamma = \partial \Omega \times (T_1, T_2)$ and $ Y_{s-1/2} (\Gamma) $ coincides with the Bochner space $L^2 ([T_1, T_2], H^{s-1/2} (\partial \Omega))$ for each 
$\alpha$ with $|\alpha|\leq s-1$.

\begin{lem} \label{l.trace.c0} Let $G$ be a relatively compact domain 
in ${\mathbb R}^n \times {\mathcal I}$ satisfying $(A1)$, $(A2)$. Then any vector 
$w \in {H}^{2m,1} (G)$  
has well-defined trace in $\partial G$. 
Besides, its derivatives  $\partial^\alpha v$ have traces 
on $\Gamma$ of the class $Y_{2m-|\alpha|-1/2} (\Gamma)$ 
for each $\alpha \in {\mathbb Z}^n_+$ with $|\alpha|\leq 2m-1$. 
\end{lem}

\begin{proof} First, we note that the space ${H}^{2m,1} (G)$  is continuously 
embedded into the isotropic Sobolev space ${H}^{1} (G)$ and hence the trace
$w_{|\partial G}$ is well-defined. Actually, it belongs to ${H}^{1/2} (\partial G)$, 
see, for instance, \cite{BeINi}. 

Next, by the structure of the domain and the Fubini theorem, 
we have 
$$
\sum_{|\alpha| \leq 2m}\|\partial ^\alpha_y w\|^2_{L^2 (G)} = 
\sum_{|\alpha| \leq 2m} 
\int_{T_1}^{T_2} \int_{G (\tau)} |\partial ^\alpha_y v(w,\tau)|^2 dy  d\tau = 
$$
$$ 
\int_{T_1}^{T_2} \| w(\cdot,\tau)\|^2 _{H^{2m} (G(\tau))}  d\tau 
$$
for any $w\in {C}^{2m} (\overline G)$. 

Now the standard Trace  Theorem for the Sobolev spaces applied for the spaces 
$H^{2m} (G(t))$ yields
$$
\|\partial ^\alpha_y h \|_{H^{2m-|\alpha|-1/2} 
(\partial G (t))}\leq C_{G(t),\alpha} \|\partial ^\alpha_y h \|_{H^{2m-|\alpha|} 
(G(t))}
$$
for all $|\alpha|\leq 2m-1$, for all $h \in H^{2m } (G (t))$ and each $t \in (T_1,T_2)$. 
Hence 
\begin{equation} \label{eq.1}
\int_{T_1}^{T_2} \sum_{|\alpha| \leq 2m-1}\|\partial ^\alpha_y w (\cdot ,\tau)\|_{H^{2m-|\alpha|-1/2} 
(\partial G (t))} d\tau \leq C(G) \int_{T_1}^{T_2} \| w(\cdot,\tau)\|^2 _{H^{2m} (G(\tau))}  
d\tau 
\end{equation}
for any $w\in {C}^{2m} (\overline G)$; 
here the positive constant 
$$
C(G)=\sup_{|\alpha|\leq 2m-1}\sup_{t\in [T_1,T_2]} 
C_{G (t),\alpha}
$$ 
is finite because the related constants $C_{G(t),\alpha} $ may be chosen to depend on 
the $(n-1)$-measure of the domains $\partial G (t)$.

According to 
\cite[\S 14]{BeINi}, any vector function $w\in {H}^{2m,1} (G)$ may be approximated 
in by vectors from ${C}_{ {\rm comp}}^\infty  ({\mathbb R}^{n+1})$ in the topology 
of the space ${H}^{2m,1} (G)$. Pick a sequence $w^{(s)}\subset {C}
_{{\rm comp}}^\infty  ({\mathbb R}^{n+1})$ approximating $w$ in ${H}^{2m,1} (G)$. 
Then \eqref{eq.1} yields that the sequence $\{ \partial ^\alpha _y w^{(s)} \}$ is 
fundamental in ${Y}_{2m-|\alpha|-1/2} (\Gamma)$ for each $\alpha$ with 
$|\alpha|\leq 2m-1$. As ${Y}_{2m-|\alpha|-1/2} (\Gamma)$ is complete we conclude 
that for each $\alpha$ with $|\alpha|\leq 2m-1$ it converges to an element $w_\alpha$, i.e. 
there is a well-defined trace $\partial ^
\alpha _y w_{|\Gamma} = w_\alpha \in {Y}_{2m-|\alpha|-1/2} (\Gamma)$ of the function 
$\partial ^\alpha _y w$ on the surface $\Gamma$. 
\end{proof}


Now we formulate the main result of this section.

\begin{thm}
\label{t.dense.base}
Let ${\mathcal L}$ satisfy assumptions $(\alpha_1)$, $(\alpha_2)$, $(\alpha_3)$, 
$(\alpha_4)$, $(UCP)$. Let  $G_1 \subset G_2 $ be  domains in ${\mathbb R}^{n}
\times \mathcal{I}$ such that bounded domain $G_1$  satisfy $(A1)$, $(A2)$.  Then 
$S_{\mathcal L}(\overline G_2)$ is everywhere dense in $\mathbf{L}^{2} _{k,\mathcal L}(G_1)$ 
if and only if for each $t \in {\mathcal I}$ the complement of the set  $G_1 (t) $   has no 
compact (non-empty) components in the set $G_2 (t)$.  
\end{thm}


\begin{proof} The proof is similar to the proof of Theorem \ref{t.dense.uniform}. However 
we need a more regular boundary for the domain $G_1$. 

\textit{Sufficiency.} Clearly, the set 
 $S_{\mathcal L}(\overline G_2)$ is everywhere dense in
 $\mathbf{L}^{2} _{k,\mathcal L}(G_1)$ if and only if the 
following relations 
\begin{equation} \label{eq.ort}
(u,w)_{\mathbf{L}^2 (G_1)} = 0 \mbox{ for all } w\in 
S_{\mathcal L}(\overline {G_2})
\end{equation}
means precisely for  the  vector $u\in \mathbf{L}^{2} _{k,\mathcal L}(G_1)$ 
that $u\equiv 0$ in $G_1$. 

Assume that the complement for each $t \in\mathcal{I}$   
the complement of the set  $G_1 (t) $   has no compact 
(non-empty) components in the set $G_2 (t)$. In order to prove the sufficiency of the statement we will again use the fact that the operator ${\mathcal L}$ admits the bilateral 
fundamental solution possessing the normality property \eqref{eq.left}. 

Let for the vector $u\in \mathbf{L}^{2} _{\mathcal L}(G_1)$ relation 
\eqref{eq.ort} is fulfilled. Consider an auxiliary vector function
\begin{equation}
\label{eq.v}
v (y,\tau) = \int_{G_1} \Phi^*(x-y,t-\tau) u (x,t)   dx \, dt.  
\end{equation}
According to \eqref{eq.right} we have
${\mathcal L}_{x,t} \Phi(x-y,t-\tau) =0 $, if  $ (x,t) \ne (y,\tau)$. 
That is why 
${\mathcal L}_{x,t} \Phi(x-y,t-\tau) =0$ in $G_2$ 
for each fixed pair $(y,\tau) \not \in \overline G_2$, i.e. 
$\Phi(x-y,t-\tau) \in S_{\mathcal L}(\overline G_2)$ 
with respect to variables $(x,t)$ for each fixed 
pair  $(y,\tau) \not \in \overline G_2$. In particular, relations  \eqref{eq.ort}  imply 
\begin{equation} \label{eq.v1}
v (y,\tau) = 0 \mbox{ in } \Big({\mathbb R}^{n} \times \mathcal{I}\Big) \setminus \overline G_2. 
\end{equation}
On the other hand, \eqref{eq.left} yields  
\begin{equation} \label{eq.v0}
{\mathcal L}^*v    = \chi_{G_1} u \mbox{ in } {\mathbb R}^{n} \times \mathcal{I},
\end{equation}
where $\chi_{ G_1}$ is the characteristic function 
of the domain  $G_1$. Obviously, 
$$
{\mathcal L}^*_{y,\tau} v (y,\tau)=0 \mbox{ in  } \Big({\mathbb R}^{n} \times \mathcal{I}\Big) 
\setminus \overline G_1,
$$
and then, by the discussed above  properties of the fundamental solution, the vector function 
 $v$ is $C^{2m,1}_{\rm loc}$-vector function in 
$\Big({\mathbb R}^{n} \times \mathcal{I}\Big)  \setminus \overline G_1 $. 

Since  the complement $G_2 (t) \setminus  G_1(t)$ has no compact components in 
$G_2(t)$  then each component of ${\mathbb R}^{n} \setminus G_1(t) $ intersects with 
${\mathbb R}^{n} \setminus  G_2(t) $ by a non-empty open set for all $t \in (T_1,T_2)$. 
Thus, \eqref{eq.v1} and the Unique Continuation Property $(UCP)$
imply that 
\begin{equation} \label{eq.v2}
v (y,\tau) = 0   \mbox{ in } \Big({\mathbb R}^{n} \times \mathcal{I}\Big)  \setminus \overline G_1 .
\end{equation}
 Then, in addition,  \eqref{eq.v0}, \eqref{eq.v2} mean that the vector  
 $v$ is a solution to the Cauchy problem 
$$
\left\{
\begin{array}{lll}
{\mathcal L}^*  v = \chi_{G_1} u \mbox{ in } {\mathbb R}^{n}  \times 
(T_1-\delta,T_2+\delta),\\
 v(y,T_2+\delta) = 0 \mbox{ on } {\mathbb R}^{n} \\
\end{array}
\right.
$$
with a sufficiently small $\delta>0$ such that $[T_1-\delta,T_2+\delta] \subset \mathcal{I}$. 
Taking in account the natural relation between parabolic and backward-parabolic operators and 
using arguments from  \cite[ch. 2, \S 5 theorem 3]{Kry08}, we may conclude that  
$v\in \mathbf{H}^{2m,1} _k ({\mathbb R}^{n} \times (T_1-\delta,T_2+\delta))$ and the 
solution is unique in this class. The regularity of this unique solution to the Cauchy 
problem can be expressed in term of the Bochner classes, too. Namely, 
 $v\in C ( [T_1-\delta,T_2+\delta], \mathbf{H}^{m}_k({\mathbb R}^{n}) )
\cap L^2 ([T_1-\delta,T_2+\delta], \mathbf{H}^{2m}_k ({\mathbb R}^{n}) )$, 
see, for instance, \cite{Lion69}, \cite[ch. 3, \S 1]{Tema79}, where similar linear problems 
for  are considered. In particular, the vector function 
 $v$ belongs to the space 
\begin{equation} \label{eq.space.1}
 C ( [T_1-\delta,T_2+\delta], \mathbf{H}^{m}_k({\mathbb R}^{n}) )
%\cap L^2 ([T_1-\delta, T_2+\delta], \mathbf{H}^{2m}_k ({\mathbb R}^{n}) )
\cap \textbf{H}^{2m,1}_k ({\mathbb R}^{n}  \times (T_1-\delta,T_2+\delta))  .
\end{equation}

\begin{lem} \label{l.dense.c0} Any vector $v$ of type \eqref{eq.v}, satisfying 
\eqref{eq.v2}, can be approximated by the vectors from 
$\mathbf{C}^\infty_{k,{\rm comp}}  (G_1)$ in the topology of the Hilbert space 
$\mathbf{H}^{2m,1}_k (G_1)$. 
\end{lem}

\begin{proof} The approximability of Sobolev functions with compact support in a 
domain is closely related to the notion of trace. 
  
Of course, as  $\textbf{H}^{2m,1}_k ({\mathbb R}^{n}  \times (T_1-\delta,T_2+\delta)) $
is embedded continuously into  $\textbf{H}^1_k (G_1)$, we see that $v \in \textbf{H}^1_{k,0}
(G_1)$, (i.e. the trace $v_{|\partial G}) \in \textbf{H}^{1/2}_{k}(\partial G_1)$ in the 
sense of isotropic Sobolev spaces and it equals to zero on $\partial G$) and then it can be 
approximated by the vectors from $\mathbf{C}^\infty_{k,{\rm comp}}  (G_1)$ in the topology of 
the Hilbert space $\textbf{H}^1_k (G_1)$. However, it is not we actually need. 

On the other hand, as the boundary of the domain $G_1$ is rather regular, it satisfies 
\cite[Condition A]{Alb} and we may use results by \cite{Alb} to conclude that $v$ may be 
approximated in $\mathbf{H}^{2m,1}_k (G_1)$ by vectors from the space 
${\mathbf C}_{k, {\rm comp}}^\infty  (\overline G_1)$. Again, it is not precisely 
what we want. 

However, \eqref{eq.space.1} implies that  traces $v_{|\overline{G(T_j)}}$ of $v$ on the 
closures of  $G(T_j)$ belong to $\mathbf{H}^{m}_k(G(T_j)) $, $1\leq j \leq 2$,  and they are 
equal to zero because of \eqref{eq.v2}. Moreover, according to Lemma \ref{l.trace.c0}, the 
traces of $\partial ^\alpha v_{|\Gamma} \in {\mathbf Y}_k^{2m-|\alpha|-1/2} (\Gamma)$ equal 
to zero on $\Gamma$ for all $|\alpha|\leq 2m-1$.

Now the statement of the lemma 
easily follows with the use of the standard regularisation, see, for instance,  
\cite[ch. 3, \S 5, \S 7,]{MikhX} for the isotropic Sobolev spaces, cf.  
see \cite{ShHeat} for cylinder domains. 

Indeed, denote by $h^{(n)}_\delta (x)$ the standard compactly supported function with 
the support in the  ball  $B(x,\delta) \subset {\mathbb R}^n$ with the centre  
at the point  $x$ and of the radius   $\delta>0$:
$$
h^{(n)}_\delta (x) = \left\{
\begin{array}{lll}
0, & \mathrm{if} & |x|\geq \delta, \\
c(\delta) \exp{(1/(|x|^2-\delta^2))}, & \mathrm{if} & |x|<\delta, \\
\end{array}
\right.
$$
where $c(\delta) $ is  the constant providing equality 
$$
\int_{{\mathbb R}^n} h_{\delta}^{(n)}(x) \, dx =1.
$$
Then, as it is well known, for any function $w \in L^1 (K)$ on a measurable compact 
$K$ in ${\mathbb R}^n$, the standard regularisation  
$$
(R ^{(n)}_\delta w) (x) = \int_{{\mathbb R}^{n}} h^{(n)}_\delta (x-y,t-\tau) w (y,\tau) \, dy d\tau 
$$
belongs to the space  $C^\infty_0  ({\mathbb R}^{n})$ for any positive number 
$\delta$ and the support of $R ^{(n)}_\delta w) $ lies in a $\delta$-neighbourhood of $K$.    

Next,  according t assumptions $(A1)$, $(A2)$ there is a real valued function 
$\rho (x,t) $ of the class $C^{2m,1}$ in a neighbourhood $U$ of the surface $\Gamma$ 
and such that 
$$
\partial G_1 (t) = \{ (x,t) \in {\mathbb R}^n \times [T_1, T_2]: 
\, \rho (x,t)<0\}, \,\, \nabla \rho \ne 0 \mbox{ in } U 
\} .
$$
Hence, for all sufficiently small numbers $\varepsilon>0$ the sets 
$$
G_1^{\varepsilon} = 
\{ (x,t) \in {\mathbb R}^n: \, \rho (x,t)<-\varepsilon\}
$$
are domains with boundaries of class $C^{2m,1}$ and 
$$
G_1 ^{\varepsilon}\Subset G_1 ^{\varepsilon'}\Subset 
G_1, 
$$
if $0<\varepsilon' <\varepsilon$, and 
$$
\lim_{\varepsilon \to +0}\mbox{mes} (G_1 \setminus 
\overline{G_1  ^{\varepsilon }}) = 0.
$$
Moreover the sets $G_1 (t)$ are domains with boundaries of class $C^{2m}$ and  for the 
Lebesgue measure of the domain  $G_1 (t) \setminus \overline{G_1 ^{\varepsilon} (t) }$ 
we have  
$$
G_1 ^{\varepsilon} (t) \Subset G_1 ^{\varepsilon'} (t)\Subset 
G_1 (t), \, t \in [T_1,T_2],  
$$
$$
\lim_{\varepsilon \to +0}\mbox{mes} (G_1 (t)\setminus 
\overline{G_1  ^{\varepsilon} (t) }) = 0
$$
uniformly with respect to $t\in [T_1,T_2]$. According to \cite[ch. 3, \S 5, lemma 1]{MikhX}, 
there is a positive constant  $C _1 (G_1)$, depending on the square of the surfaces 
$G_1 (T_j) $, $1\leq j\leq 2$, only, and such that 
\begin{equation} \label{eq.est.12}
C_1 (G_1)\, \varepsilon ^{-2}
 \Big(
\int_{T_1}^{T_1+\varepsilon}\| \tilde v \|^2_{L^2 (G (t) )} dt +
\int_{T_2-\varepsilon}^{T_2}\| \tilde v \|^2_{L^2 (G (t) )} dt 
\Big)\leq
\end{equation}
\begin{equation*}  
\int_{T_1}^{T_1+\varepsilon}(\| \tilde v \|^2_{L^2 (G (t) )} +
\| \partial_t \tilde v \|^2_{L^2 (G (t) )} )dt +
\int_{T_2-\varepsilon}^{T_2}(\| \tilde v \|^2_{L^2 (G (t) )} +
\| \partial_t \tilde v \|^2_{L^2 (G (t) )}) dt 
\end{equation*}
for any function $\tilde v \in H^{1} (G_1)$ with zero traces $\tilde v_{|G_1 (T_j)}$
on $G_1 (T_j)$, $1\leq j\leq 2$. 

Similarly, under assumptions $(A1)$, $(A2)$, if  $\partial G(t) \in C^{2m}$ then 
there is a constant  $C _0 (\Gamma)$, depending on the square of the surface 
$ \Gamma$, only, and such that  
\begin{equation} \label{eq.est.11}
\Big(\int_{T_1}^{T_2}\|\partial ^\alpha \tilde v \|^2_{L^2 (G (t) \setminus \overline {G^
\varepsilon}(t))} dt \Big)^{1/2}\leq C_0 (\Gamma)\, \varepsilon  ^{2m-|\alpha|} 
\Big( \int_{T_1}^{T_2} \|\tilde v \|^2_{H^{2m}(G (t) \setminus \overline {G^\varepsilon}(t))} 
dt \Big)^{1/2}
\end{equation}
for any function $\tilde v \in H^{2m,1} (G_1)$ with zero traces 
$\partial^{\alpha}\tilde v_{|\Gamma}$ on $\Gamma$ corresponding to $|\alpha|\leq 2m-1$.

Set 
$$
R_\varepsilon (x,t) = \int_{T_1+\varepsilon/5}^{T_2-\varepsilon/5} 
 h^{(1)}_{\varepsilon/7} (t-\tau) \, \int_{G_1^{\varepsilon/5} (\tau)}  h^{(n)}_{\varepsilon/7} (x-y) \, dy \, d\tau.
$$
Using assumptions 
on $G_1$ we may choose some constants $c_{\alpha}$, $c_{i} $ independent on  $x$ and $t$ 
such that 
\begin{equation} \label{eq.Bochner.1}
0\leq R_\varepsilon  (x,t)\leq 1,
\end{equation}
\begin{equation} \label{eq.Bochner.2}
|\partial^\alpha_x R_\varepsilon (x,t)| \leq c_{\alpha} \, \varepsilon ^{-|\alpha|}, 
\,\, 
|\partial^i_t  R_\varepsilon (x,t)| \leq c_{i} \, \varepsilon ^{-i}, 
\end{equation}
for all $x \in {\mathbb R}^n$, $t \in \mathbb R$, 
$\alpha \in {\mathbb Z}^n_+$, $i \in {\mathbb Z}_+$, 
cf., for instance, \cite[ch. 3, \S 5]{MikhX}. 

Fix a sequence $\{ v_j\} \subset C^\infty _0({\mathbb R}^{n+1})$, converging to 
$v$ in the space ${\mathbf H}_k^{2m,1} ({\mathbb R}^{n+1} )$ (as we mentioned above, 
the existence of such a sequence follows from \cite[\S 14]{BeINi}). 
Then the functional sequence 
$$
 \{  v_{j,\varepsilon} (x,t)=  R_\varepsilon (x,t)v_j (x,t)\}
$$
lies to ${\mathbf C}^\infty _{k,0} (G_1)$. 

By the triangle inequality, 
\begin{equation} \label{eq.vk.1}
\|v - v_{j,\varepsilon} \|_{{\mathbf H}_k^{2m,1} (G_1)} \leq 
\|v - v_{j} \|_{{\mathbf H}_k^{2m,1} (G_1)} + 
\|v _k - v_{j,\varepsilon} \|_{{\mathbf H}_k^{2m,1} (G_1)} .
\end{equation}
As 
\begin{equation} \label{eq.lim.1}
\lim_{j\to + \infty}\|v - v_{j} \|_{{\mathbf H}_k^{2m,1} (G_1)} =0,
\end{equation}
then we need to estimate the second summand in the right hand side of formula 
\eqref{eq.vk.1}, only. However,  
$$
v _j (x,t)- v_{j,\varepsilon} (x,t) = 
(1-R_\varepsilon (x,t)) v_j (x,t)
$$
and, in particular,  
\begin{equation} \label{eq.eps.0}
v _j (x,t)- v_{j,\varepsilon} (x,t) = 0 \mbox{ for all } (x,t) \in 
\cup_{t=T_1+\varepsilon}^{T_2-\varepsilon}G_1 ^{\varepsilon} (t) .
\end{equation}
Hence, 
\begin{equation} \label{eq.eps}
2^{-1}\|v _j - v_{j,\varepsilon}\|^2_{H^{2m,1} (G_1)} \leq  
\|(1-R_\varepsilon)  \partial_t v_j \|^2
_{L^{2} (G_1)} + 
\end{equation}
$$
\Big\|  (\partial_t R_\varepsilon) 
v_j \Big\|^2 _{L^{2} (\cup_{t\in [T_1,T_1+\varepsilon] \cup [T_2-\varepsilon ,T_2]}
G_1 ^{\varepsilon} (t) )} +
$$
$$
\sum_{ |\alpha| \leq 2m}
\| \partial^{\alpha} ((1-R_\varepsilon)   v_j) \|^2
_{L^{2} (\cup_{t\in [T_1,T_2]} G_1 (t)  \setminus G_1^{\varepsilon} (t))} . 
$$
Then \eqref{eq.est.12}, 
\eqref{eq.est.11},  \eqref{eq.Bochner.1}, \eqref{eq.Bochner.2}, \eqref{eq.eps.0} 
and the Fubini theorem imply that  
\begin{equation} \label{eq.eps.1}
\sum_{ |\alpha| \leq 2m}
\| \partial^{\alpha} (1-R_\varepsilon   v_j) \|^2
_{L^{2} (G_1)}  
\leq 
\end{equation}
$$
\sum_{|\alpha| \leq 2m } 
\|(1-R_\varepsilon)(  \partial^{\alpha} v_j \|^2
_{L^{2} (\cup_{t\in [T_1,T_2]} G_1 (t)  \setminus G_1^{\varepsilon} (t))}  +
$$
$$
\sum_{|\beta+\gamma| \leq 2m , \beta \ne 0} 
\|(\partial^{\beta}R_\varepsilon)(  \partial^{\gamma} v_j \|^2
_{L^{2} (\cup_{t\in [T_1,T_2]} G_1 (t)  \setminus G_1^{\varepsilon} (t))}  \leq 
$$
$$
\sum_{|\alpha| \leq 2m } 
\|(1-R_\varepsilon)(  \partial^{\alpha} v_j \|^2
_{L^{2} (\cup_{t\in [T_1,T_2]} G_1 (t)  \setminus G_1^{\varepsilon} (t))}  +
$$
$$
\sum_{|\beta+\gamma| \leq 2m , \beta \ne 0} 
\varepsilon^{2m-|\gamma+\beta|}\|   v_j \|^2
_{H^{2m,0} (\cup_{t\in [T_1,T_2]} G_1 (t)  \setminus G_1^{\varepsilon} (t))}  \leq 
$$
$$
\| v\|_{H^{2m,1} (\cup_{t\in [T_1,T_2]}  
G_1 (t) \setminus \overline{G_1^{\varepsilon} (t) })}
$$
with constants $C,\tilde C$, independent on $v$ and $\varepsilon$. 

The boundaries of the domains $\cup_{t \in [T_1,T_1+\varepsilon] } G_1 (t)$ and 
$\cup_{t \in [T_2-\varepsilon,T_2] } G_1 (t)$ are not smooth, but combining results 
\cite[ch. 3, \S 5]{MikhX} related to a function $v$, having the trace vanishing on surfaces  
$G_1(T_1)$ and $G_2 (T_2)$, with bounds  \eqref{eq.Bochner.1}, \eqref{eq.Bochner.2}, 
we see that 
\begin{equation} \label{eq.eps.3}
\| (\partial _t R_\varepsilon) v \|^2
_{L^{2} (\cup_{t\in [T_1,T_1+\varepsilon] \cup [T_2-\varepsilon ,T_2]}
G_1 ^{\varepsilon} (t) )}
\leq 
\end{equation}
$$
\varepsilon^{-1} \cdot \varepsilon 
C \, \Big( \| _t v\|^2
_{L^{2} (\cup_{t\in [T_1,T_1+\varepsilon] \cup [T_2-\varepsilon ,T_2]} 
G_1 ^{\varepsilon} (t) )} + 
\| \partial _t v\|^2
_{L^{2} (\cup_{t\in [T_1,T_1+\varepsilon] \cup [T_2-\varepsilon ,T_2]} 
G_1 ^{\varepsilon} (t) )} \Big)\leq 
$$
$$
\| v\|_{H^{2m,1} (\cup_{t\in [T_1,T_1+\varepsilon] \cup [T_2-\varepsilon,T_1]}  
G_1 (t) \setminus \overline{G_1^{\varepsilon} (t) })}
$$
with a constant $C$, independent on $v$ and  $\varepsilon$.

Besides, according to  \eqref{eq.Bochner.1}, \eqref{eq.Bochner.2}, 
\begin{equation} \label{eq.eps.4}
\|(1-R_\varepsilon)\partial^i_t v \|^2
_{L^{2} (G_1)}
\leq C \| v\|_{H^{2m,1} (\cup_{t\in [T_1,T_2]}  
G_1 (t) \setminus \overline{G_1^{\varepsilon} (t) })}
\end{equation}
with a constant $C$, independent on $v$ and  $\varepsilon$.

Using the continuity of the Lebesgue integral with respect 
to the measure of the integration set, we conclude that 
\begin{equation} \label{eq.lim.2}
\lim_{\varepsilon \to + 0}
\| v\|_{H^{2m,1} (\cup_{t\in [T_1,T_1+\varepsilon] \cup [T_2-\varepsilon,T_1]}  
G_1 (t) )} =0,
\end{equation}
\begin{equation} \label{eq.lim.2A}
\lim_{\varepsilon \to + 0}\| v\|_{H^{2m,1} (\cup_{t\in [T_1,T_2]}  
G_1 (t) \setminus \overline{G_1^{\varepsilon} (t) })} =0.
\end{equation}

Finally, combining estimates \eqref{eq.eps}--\eqref{eq.eps.4} and taking in account 
\eqref{eq.lim.2}, \eqref{eq.lim.2A}, 
we conclude that the statement of the lemma holds true. 
\end{proof}

Next, using lemma \ref{l.dense.c0} and fixing a sequence 
$\{ v_k \} \subset \mathbf{C}^\infty _{k, {\rm comp}} 
(G_1)$ converging to the vector 
$v$ in $\mathbf{H}^{2m,1}_k (G_1)$, we see that  
$$
\|u\|^2_{\mathbf{L}^2_k (G_1)} = 
(u , {\mathcal L}^* v )_{\mathbf{L}^2_k (G_1)} = 
\lim_{i \to + \infty }(u, {\mathcal L}^* v_ i )_{\mathbf{L}^2 _k(G_1)}   
= 0,
$$
because ${\mathcal L} u =0 $ in $G_1$ in the sense of distributions. 
Thus,  $u \equiv 0$ in $G_1$, i.e. the sufficiency is proved.

\textit{Necessity.} Actually, we may use the arguments similar to 
that in the proof of the necessity of Theorem \ref{t.dense.uniform}. 
Indeed, let there be a number $t_0\in \mathbb R$ such that the set  $G_1 (t_0) $   
have a compact (non-empty) component $K_0$ in the set $G_2 (t_0)$. By the assumptions on the 
boundaries of the set $G_1 (t)$, we may consider $K_0$ as the closure of a domain $D_0$ with 
Lipschitz boundary in ${\mathbb R}^n$. Fix a point $y_0 \in D_0$. Then any vector column 
$U_j (x,t)$, $1\leq j \leq k$, of the fundamental matrix $\Phi (x-y_0, t-t_0)$ belongs to the 
space ${\mathbf L}^2_{k,\mathcal L} (G_1)$. If the vector function $U_j (x,t)$ can be 
approximated in ${\mathbf L}^2_k  (G_1)$ by a sequence $\{ u^{(i)} \}_{i \mathbb N}$ from the 
space $S_{\mathcal L} (\overline G_2)$ then this sequence should satisfy the Cauchy criterion 
in the Banach space ${\mathbf L}^2_k  (G_1 \setminus G_3) $ with a neighbourhood $G_3$ of the 
point $(y_0, t_0)$. Then by the standard a priori estimates for parabolic systems we conclude 
that sequence $\{ u^{(i)} \}_{i \mathbb N}$ converges, actually, to a vector function 
$U(x,t)$ in the space of smooth $k$-vector functions over a neighbourhood of the point 
$(y_0,t_0)$. In particular, assumption $(UCP)$ to the system ${\mathcal L}$ with respect to 
the space variables implies that the vector function $U_j (x,t) \in S_{\mathcal L} 
(({\mathbb R}^n \times (0,+\infty)) \setminus (y_0,t_0))$ uniquely extends to a neighbourhood 
$G_4$ of the point $(y_0,t_0)$ as a solution $\tilde U_j (x,t)$ to ${\mathcal L}$. Thus, we 
obtain a contradiction because for the matrix $\tilde \Phi (x,t)$ with columns $\tilde U_j 
(x,t)$, $1\leq j \leq k$, we have ${\mathcal L} \tilde \Phi (x,t) =0$ in $G_4$ but
${\mathcal L} \Phi  (x-y_0,t-t_0)$ coincides with the $\delta$-functional concentrated at 
the point $(y_0,t_0)$. 
\end{proof}

Next, we obtain the following useful statement.

\begin{cor} \label{c.dense.base0} 
Let the coefficients of the operator $\mathcal L$ be smooth, bounded and real analytic with 
respect to the variables $x$ for each $t\in \mathcal{I}$. Let also $s,\gamma\in 
{\mathbb Z}_+$,  $G_1 \subset G_2 $ be domains in ${\mathbb R}^{n}\times \mathcal{I}$ such that 
bounded $G_1$ satisfy $(A1)$, $(A2)$. If  for each $t \in \mathcal{I}$ the complement 
of the set  $G_1 (t) $ has no compact (non-empty) components in the set $G_2 (t)$ then the 
spaces ${\mathbf C}^\infty_{k,\mathcal L}(\overline G_2)$ and ${\mathbf H}^{\gamma,2s,s} 
_{k,\mathcal L}(G_2)$ are everywhere dense in  ${\mathbf L}^{2} _{k,\mathcal L}(G_1)$. 
\end{cor}

\begin{proof} Follows immediately from theorem \ref{t.dense.base}, 
because of embeddings \eqref{eq.emb.H}.
\end{proof}

Finally, Theorem \ref{t.dense.base} allows to prove the existence 
of a basis with the double orthogonality property in the spaces of solutions 
to the operator $\mathcal L$ that is very useful to investigate the non-standard 
ill-posed Cauchy for elliptic and parabolic equations, see, 
\cite[Ch. 12]{Tark36}, \cite{ShTaLMS}, \cite{VKShLame}.

\begin{cor} \label{c.bdo} 
Let ${\mathcal L}$ satisfy assumptions $(\alpha_1)$, $(\alpha_2)$, $(\alpha_3)$, 
$(\alpha_4)$, $(UCP)$. Let  $G_1 \subset G_2 $ be domains in ${\mathbb R}^{n}
\times \mathcal{I}$, such that $T_1 (G_1)= T_1 (G_2)$, $T_2 (G_1) = T_2 (G_2)$ and 
 bounded domain $G_1$ satisfy $(A1)$, $(A2)$. If for each $t \in \mathcal{I}$ the complement of the set 
$G_1 (t) $ has no compact (non-empty) components in the set $G_2 (t)$ then there is an 
orthonormal basis  $\{ b_\nu\}$ is the space $\mathbf{H}^{2m,1}_{k,\mathcal L} (G_2)$ such 
that its restriction $\{ b_{\nu|G_1}\}$ to $G_1$ is on orthonormal basis in the space 
$\mathbf{L}^{2}_{k,\mathcal L} (G_1)$. 
\end{cor} 

\begin{proof} By the definition, the space $\mathbf{H}^{2m,1} _{k,\mathcal L} (G_2)$  is 
embedded continuously into the space $\mathbf{L}^2 _{k,\mathcal L} (G_1)$.  We denote by $R$ 
the natural embedding operator  
$$
R: \mathbf{H}^{2m,1} _{k,\mathcal L}(G_2)\to 
\mathbf{L}^2 _{k,\mathcal L}(G_1).
$$ 
As the numbers $T_1$ and $T_2$ are the same for the domains $G_1$ and $G_2$, 
the Unique Continuation Property $(UCP)$ for the operator  $\mathcal L$ 
with respect to the space variables implies that the operator $R$ is injective.
Besides, it follows from Theorem \ref{t.dense.base} that the range of the operator 
$R$ is everywhere dense in the space ${\mathbf L}^2 _{k,\mathcal L}(G_1)$. 

By the definition, the anisotropic Sobolev space   
$\mathbf{H}^{2m,1} _{k,\mathcal L}(G_2)$ is embedded continuously 
to the isotropic Sobolev space  $\mathbf{H}^{1} _{k,\mathcal L}(G_2)$. Besides,  
by Rellich-Kondrashov theorem the embedding  
$\mathbf{H}^{1} _{k}(G_1) \to {\mathbf L }^2_k (G_1)$ is compact. 
Thus, taking in account the continuous embedding 
$\mathbf{H}^{1} _{k}(G_2) \to {\mathbf H }^1_k (G_1)$, we see that 
the natural embedding operator $R$  is compact. 

Finally,   \cite[example 1.9]{ShTaLMS} implies that the complete system of eigen-vectors 
of the  compact self-adjoint operator
 $R^* R: 
\mathbf{H}^{2m,1} _{k,\mathcal L}(G_2) \to \mathbf{H}^{2m,1} _{k,\mathcal L}(G_1)$ 
is the basis looked for; here  $R^*$ is the adjoint operator for $R$ 
in the sense of the Hilbert space theory.
\end{proof}

\begin{cor} \label{c.bdo.s} 
Let the coefficients of the operator $\mathcal L$ be smooth, bounded and real analytic with 
respect to the variables $x$ for each $t\in \mathcal{I}$. Let also $s,\gamma\in 
{\mathbb Z}_+$,  $G_1 \subset G_2 $ be domains ${\mathbb R}^{n}
\times \mathcal{I}$, such that $T_1 (G_1)= T_1 (G_2)$, $T_2 (G_1) = T_2 (G_2)$ and 
 bounded  domain $G_1$ satisfy $(A1)$, $(A2)$.  If for each $t \in \mathcal{I}$ the complement of the set  
$G_1 (t) $ has no compact (non-empty) components in the set $G_2 (t)$ then there is an 
orthonormal basis $\{ b_\nu\}$ is the space $\mathbf{H}^{\gamma,2ms,s}_{k,\mathcal L} (G_2)$ 
such that its restriction $\{ b_{\nu|G_1}\}$ to $G_1$ is on orthonormal basis in the space  
$\mathbf{L}^{2} _{k,\mathcal L} (G_1)$. 
\end{cor} 

\bigskip

{\sc Acknowledgments.}
The first author was supported 
by the Krasnoyarsk Mathematical Center and financed by the Ministry of Science and Higher 
Education of the Russian Federation (Agreement No.075-02-2023-936). The 
second author was supported by Supported by the Ministry of Science and Higher Education of 
the Russian Federation, (Agreement 075-10-2021-093, Project MTH-RND-2124).
 

\begin{thebibliography}{XXXXXX}

\bibitem{Alb}
Alborova M.S. \textit{A density theorem}, Vladikavkaz Mathematical Journal, 
{\bf 1} (2001), no.~3, 1--7. 
 
\bibitem{Beheshti2016} 
{Beheshti M, Umapathy K, Krishnan S.}
\textit{Electrophysiological Cardiac Modeling: A Review}.
Crit Rev Biomed Eng. 2016; 44(1--2):99--122.

\bibitem{BeINi} Besov O.V., Il'in V.P., Nikol'skii S.M., Integral representation 
of functions and imbedding  theorems, V.1.
Washington, DC, V.H. Winston $\&$ Sons, 1978.

\bibitem{D80}
{ Diaz, R.},
  \textit{ A Runge theorem for solutions of the heat equation},
  Proc. Amer. Math. Soc. {\bf 80} (1980), no.~4, 643--646. 

\bibitem{Edw65} Edwards R.E. {Functional anlysis}, Holt, Rinehart and Winston, 
London, 1965.

\bibitem{eid} Eidel'man S.D., 
\textit{Parabolic equations}, 
Partial differential equations -- 6, Itogi Nauki i Tekhniki. Ser. Sovrem. Probl. Mat. Fund. 
Napr., 63, VINITI, Moscow, 1990, 201--313.

\bibitem{frid}
{Friedman, A.}, {Partial differential equations of parabolic type}, 
Englewood Cliffs, NJ, Prentice-Hall, Inc., 1964.

\bibitem{GauTa10}
{Gauthier, P.M., Tarkhanov, N,} 
\textit{Rational approximation and universality for a quasilinear parabolic equation}, 
Journal of Contemporary Mathematical Analysis, V. 43 (2008), 353--364.

\bibitem{H68}
Havin, V. P., \textit{Approximation by analytic functions in the mean}, 
Dokl. Akad. Nauk SSSR, 178:5 (1968),  1025--1028. 

\bibitem{J}
{Jones, B.F., Jr.},
  \textit{ An approximation theorem of Runge type for the heat equation},
  Proc. Amer. Math. Soc. {\bf 52} (1975), no.~1, 289--292.

\bibitem{KSU_ZAMM}
Kalinin V.,  Shlapunov A.A.,  Ushenin K.,
\textit{On uniqueness theorems for the inverse problem of Electrocardiography in the Sobolev 
spaces}, Z. Angew. Math. Mech. V. 103:1 (2023), e202100217.

\bibitem{Kry08} {Krylov, N.V.},
{Lectures on elliptic and parabolic equations in Sobolev spaces}. Graduate Studies in 
Math. V. 96,  AMS, Providence, Rhode Island, 2008.
	
\bibitem{Kry96} {Krylov, N.V.},
{Lectures on elliptic and parabolic equations in H\"older spaces}. Graduate Studies in 
Math. V. 12,  AMS, Providence, Rhode Island, 1996.

\bibitem{KuSh}
{Kurilenko, I.A., Shlapunov, A.A.}, 
\textit{On Carleman-type Formulas for Solutions to the Heat Equation}, 
Journal of Siberian Federal University, Math. and Phys.,  12:4 (2019), 421--433.

\bibitem{LadSoUr67}
{Ladyzhenskaya, O.A., Solonnikov V.A., Ural'tseva N.N.},
{Linear and quasilinear equations of parabolic type}, M., Nauka,  1967. 

\bibitem{Lion69}
{Lions, J.-L.}, {Quelques m\'{e}thodes de r\'{e}solution des probl\`{e}mes aux 
limites non lin\'{e}are}, Dunod/Gauthier-Villars, Paris, 1969. 

\bibitem{Mal56}
{Malgrange, B.},  
\textit{Existence et approximation des solutions des \'equations aux d\'eriv\'ees 
partielles et des \'equations de convolution}, Annales de l'Institut Fourier, V. 6 (1956),  271--355. 

\bibitem{Mal63}
{Malgrange, B.},  
\textit{Sur les syst\'emes differentiels \'a coefficients constants}, Colloq. internat. 
Centre nat. rech. scient., 117 (1963),  113--122. 
 
\bibitem{MaHa74}
 Maz'ya, V.G., Havin,  V.P.,  
\textit{The solutions of the Cauchy problem for the Laplace equation (uniqueness, normality, 
approximation)}, Tr. Mosk. Mat. Obs., 30 (1974),  61--114.

\bibitem{MMT17}
Makhmudov K.O., Makhmudov O.I.,  Tarkhanov N.N., 
\textit{A Nonstandard Cauchy Problem for the Heat Equation}, 
Math. Notes, 102:2 (2017), 250--260. 

\bibitem{Merg56}
 Mergelyan, S.N., 
\textit{Harmonic approximation and approximate solution of the Cauchy problem for the Laplace 
equation}, Uspekhi Mat. Nauk, 11:5(71) (1956),  3--26. 

\bibitem{MikhX}
{Mikhailov, V.P.}, 
{Partial differential equations}, M., Nauka, 1976. 

\bibitem{Pal}
Palamodov, V.P. {Linear Differential Operators with Constant}, 
Coefficients, Grund\-leh\-ren der Mathematischen Wissenschaften 168. Springer, Berlin, 1970.

\bibitem{PuSh15}
{Puzyrev, R.E., Shlapunov, A.A.} 
\textit{ On a mixed problem for the parabolic Lam\'e type operator}. 
J. Inv. Ill-posed Problems,  23:6 (2015), 555--570.

\bibitem{R1885}
{Runge, C.}, \textit{Zur Theorie der eindeutigen analytischen Funktionen}, 
Acta Math. 6 (1885), 229--244. 

\bibitem{ShHeat}
{Shlapunov, A.A.},
 \textit{On approximation of solutions to the heat 
equation from Lebesgue class $L^2$ by more regular solutions}, 
Math. Notes, {\bf 111} (2022), no.~5, 778--794.

\bibitem{ShTaLMS} 
{Shlapunov, A.A.,  Tarkhanov, N.}, 
\textit{ Bases with double orthogonality in the Cauchy problem for systems with 
injective symbols}. Proc. London. Math. Soc.,  71 (1995), N. 1, p. 1--54. 

\bibitem{Sch71} Schaefer H. {Topological vector spaces}, Springer-Verlag, Berlin, 1971.

\bibitem{sol}
Solonnikov, V.A., 
 \textit{On boundary value problems for linear parabolic systems of differential equations of 
general form}, Proc. Steklov Inst. Math., 83 (1965), 1--184.

\bibitem{Tark36}
{Tarkhanov, N.},
{The Cauchy Problem for Solutions of Elliptic Equations},
 Akademie-Verlag, Berlin, 1995.

\bibitem{Tark37}
{Tarkhanov, N.},
 {The Analysis of Solutions of Elliptic Equations},
 Kluwer Academic Publishers, Dordrecht, NL, 1997.

\bibitem{Tema79}
Temam, R.,
{Navier-Stokes Equations. Theory and Numerical Analysis},
  North Holland Publ. Comp., Amsterdam, 1979.

\bibitem{UDP} 
Uspenskii S.V., Demidenko G.V., Perepelkin V.G., 
{Embedding theorems and applications to differential equations}, 
 Nauka, Novosibirsk, 1984. 


\bibitem{VKShLame}
Vilkov P.Yu., Kurilenko I.A., Shlapunov A.A., 
\textit{Approximation of solutions to parabolic Lam\'e type operators in cylinder domains and 
Carleman's  formulas for them},  Siberian Math. J., Vol. 63, No. 6, 2022, pp. 1049--1059. 

\bibitem{V67}
Vitushkin, A.G, 
\textit{The analytic capacity of sets in problems of approximation theory}, 
Uspekhi Mat. Nauk, 22:6(138) (1967),  141--199. 

\bibitem{W1885} 
Weierstra\ss{}, K., \textit{Uber die analytische Darstellbarkeit sogenannter 
willkurlicher Fu\-nc\-tio\-nen einer reellen Veranderlichen Sitzungsberichteder}, 
Sitzungsberichte der Akademie zu Berlin, 1885, pp. 63--639 and 789--805.

\end{thebibliography}
\end{document}
